% ============================================================
%  Card (1995) — DML & Kappa Weighting Estimators
%  Results tables and findings write-up
% ============================================================
\documentclass[12pt]{article}

\usepackage[a4paper, margin=2.5cm]{geometry}
\usepackage{booktabs}
\usepackage{threeparttable}
\usepackage{siunitx}
\usepackage{amsmath}
\usepackage{amssymb}
\usepackage{caption}
\usepackage{setspace}
\usepackage{microtype}
\usepackage[round,authoryear]{natbib}
\usepackage[hidelinks]{hyperref}
\usepackage[table]{xcolor}
\usepackage{colortbl}

% ---- diagnostic colour palette ----
\definecolor{diagOK}{RGB}{234,243,222}       % light green  — Sum_w=0, ESS≥2
\definecolor{diagWarn}{RGB}{250,238,218}     % light amber  — Sum_w≠0 (not transl.-inv.)
\definecolor{diagFlag}{RGB}{252,235,235}     % light red    — ESS=1 or ESS=0
\definecolor{diagFlagText}{RGB}{163,45,45}   % dark red     — text in flagged cells
\definecolor{diagWarnText}{RGB}{133,79,11}   % dark amber   — text in warn cells
\definecolor{diagOKText}{RGB}{59,109,17}     % dark green   — text in ok cells

\onehalfspacing

\sisetup{
  table-format        = 1.3,
  table-number-alignment = center,
  detect-all
}


\onehalfspacing

\sisetup{
  table-format        = 1.3,
  table-number-alignment = center,
  detect-all
}

\newcommand{\se}[1]{{\small(#1)}}
\newcommand{\tstat}[1]{{\footnotesize[#1]}}

% ============================================================
\begin{document}

\section*{Results: Card (1995) --- Returns to Schooling}

\subsection*{Setup}

We apply the double machine learning (DML) framework of
\citet{chernozhukov2018double} and the kappa weighting estimators of
\citet{sloczynski2025abadie} to the Card (1995) returns-to-schooling
data. The instrument is \texttt{nearc4} (proximity to a four-year
college). We consider two treatment definitions and two covariate
specifications, yielding four estimation cells:

\begin{itemize}
  \item \textbf{Treatment:} \texttt{somecol} ($D_1 = \mathbf{1}\{\text{educ} > 12\}$,
        any post-secondary education) and
        \texttt{educ16} ($D_2 = \mathbf{1}\{\text{educ} \geq 16\}$,
        college degree or more).
  \item \textbf{Covariate specification:} the rich \emph{Card} specification
        (\texttt{exper}, \texttt{expersq}, \texttt{black}, \texttt{south},
        \texttt{smsa}, \texttt{smsa66}, regional dummies
        \texttt{reg661}--\texttt{reg668}) and the parsimonious
        \emph{Kitagawa} specification
        (\texttt{black}, \texttt{south}, \texttt{smsa}, \texttt{smsa66},
        \texttt{south66}).
\end{itemize}

The outcome is log wages measured in dollars ($Y_{\$}$). DML estimates
are obtained via five-fold cross-fitting with tuned random forests
and XGBoost as learners. Kappa estimates use analytic
M-estimation standard errors.

% ============================================================
\subsection*{Table 1: DML Estimators (PLR-IV and Wald-AIPW)}

\begin{table}[htbp]
  \centering
  \begin{threeparttable}
    \caption{Double machine learning IV estimates of the returns to schooling,
             Card (1995). Outcome: log wages (dollars).
             Instrument: \texttt{nearc4}.}
    \label{tab:dml_estimates}
    \begin{tabular}{l
                    S[table-format=1.3]
                    S[table-format=1.3]
                    S[table-format=1.3]
                    S[table-format=1.3]}
      \toprule
      & \multicolumn{2}{c}{\texttt{somecol} ($D_1$)}
      & \multicolumn{2}{c}{\texttt{educ16} ($D_2$)} \\
      \cmidrule(lr){2-3} \cmidrule(lr){4-5}
      Estimator
        & {Card} & {Kitagawa}
        & {Card} & {Kitagawa} \\
      \midrule
      PLR-IV
        & 0.681$^{**}$ & 0.593$^{*}$ & 1.140$^{**}$ & 1.030 \\
        & \se{0.289}   & \se{0.340}  & \se{0.535}   & \se{0.684} \\
        & \tstat{2.358}& \tstat{1.742}& \tstat{2.131}& \tstat{1.507} \\[4pt]
      Wald-AIPW
        & 0.279 & 0.292 & 0.530 & 0.535 \\
        & \se{0.214} & \se{0.249} & \se{0.398} & \se{0.462} \\
        & \tstat{1.303} & \tstat{1.173} & \tstat{1.332} & \tstat{1.157} \\
      \bottomrule
    \end{tabular}
    \begin{tablenotes}[flushleft]\footnotesize
      \item Standard errors in parentheses; $t$-statistics in brackets.
            $^{***}p<0.01$, $^{**}p<0.05$, $^{*}p<0.10$.
            DML uses five-fold cross-fitting with tuned random forests
            and XGBoost. Card specification: \texttt{exper}, \texttt{expersq},
            \texttt{black}, \texttt{south}, \texttt{smsa}, \texttt{smsa66},
            \texttt{reg661}--\texttt{reg668}.
            Kitagawa specification: \texttt{black}, \texttt{south},
            \texttt{smsa}, \texttt{smsa66}, \texttt{south66}.
    \end{tablenotes}
  \end{threeparttable}
\end{table}

% ============================================================
\subsection*{Table 2: Kappa Weighting Estimators}

\begin{table}[htbp]
  \centering
  \begin{threeparttable}
    \caption{Kappa weighting estimators of the LATE,
             Card (1995). Outcome: log wages (dollars).
             Instrument: \texttt{nearc4}.
             Analytic M-estimation standard errors.}
    \label{tab:kappa_estimates}
    \begin{tabular}{l
                    S[table-format=1.3]
                    S[table-format=1.3]
                    S[table-format=1.3]
                    S[table-format=1.3]}
      \toprule
      & \multicolumn{2}{c}{\texttt{somecol} ($D_1$)}
      & \multicolumn{2}{c}{\texttt{educ16} ($D_2$)} \\
      \cmidrule(lr){2-3} \cmidrule(lr){4-5}
      Estimator
        & {Card} & {Kitagawa}
        & {Card} & {Kitagawa} \\
      \midrule
      $\hat{\tau}^{cb}_{u}$ (CBPS)
        & 0.376$^{*}$ & 0.331 & 0.853 & 0.588 \\
        & \se{0.223}  & \se{0.236} & \se{0.549} & \se{0.433} \\[4pt]
      $\hat{\tau}^{ml}_{u}$ (MLE, no trim)
        & 0.331 & 0.356 & 0.619 & 0.628 \\
        & \se{0.202} & \se{0.244} & \se{0.387} & \se{0.448} \\[4pt]
      $\hat{\tau}^{ml}_{a,10}$ (MLE, 10\% trim)
        & 0.346$^{*}$ & 0.293 & 0.586$^{*}$ & 0.836 \\
        & \se{0.200}  & \se{0.252} & \se{0.356} & \se{0.821} \\[4pt]
      \midrule
      \multicolumn{5}{l}{\textit{Scaled / unnormalised variants}} \\[2pt]
      $\hat{\tau}^{ml}_{a}$
        & 0.170 & 0.842$^{**}$ & 0.315 & 1.617$^{*}$ \\
        & \se{0.370} & \se{0.362} & \se{0.696} & \se{0.891} \\[4pt]
      $\hat{\tau}^{ml}_{t}$
        & 0.171 & 0.769$^{**}$ & 0.319 & 1.367$^{**}$ \\
        & \se{0.367} & \se{0.308} & \se{0.687} & \se{0.648} \\[4pt]
      $\hat{\tau}^{ml}_{a,0}$
        & 0.154 & 1.066$^{*}$ & 0.266 & 2.712 \\
        & \se{0.354} & \se{0.574} & \se{0.639} & \se{2.577} \\
      \bottomrule
    \end{tabular}
    \begin{tablenotes}[flushleft]\footnotesize
      \item Standard errors in parentheses (analytic M-estimation).
            $^{***}p<0.01$, $^{**}p<0.05$, $^{*}p<0.10$.
            $\hat{\tau}^{cb}_{u}$: CBPS propensity score, untrimmed.
            $\hat{\tau}^{ml}_{u}$: MLE propensity score, untrimmed.
            $\hat{\tau}^{ml}_{a,10}$: MLE with 10\% trimming.
            $\hat{\tau}^{ml}_{a}$, $\hat{\tau}^{ml}_{t}$,
            $\hat{\tau}^{ml}_{a,0}$: alternative scalings of the
            unnormalised kappa estimator (see \citealt{sloczynskietal2025}).
    \end{tablenotes}
  \end{threeparttable}
\end{table}

% ============================================================
\subsection*{Discussion of Findings}

\paragraph{PLR-IV versus Wald-AIPW.}
Across all four specification cells, the PLR-IV estimator
consistently yields larger point estimates than Wald-AIPW,
by a factor of roughly two. Under the Card specification with
\texttt{somecol}, PLR-IV returns an estimate of $0.681$
(SE $= 0.289$, $p = 0.018$), while Wald-AIPW yields $0.279$
(SE $= 0.214$, $p = 0.193$). The same pattern holds for
\texttt{educ16}: PLR-IV gives $1.140$ (significant at 5\%)
versus $0.530$ for Wald-AIPW (insignificant). 



\paragraph{Treatment definition: \texttt{somecol} versus \texttt{educ16}.}
Estimates roughly double when the treatment switches from
any post-secondary education (\texttt{somecol}) to attaining
a college degree (\texttt{educ16}). This is expected: the
complier group for \texttt{educ16} is narrower---individuals
pushed to complete four years of college by proximity to a
college---and the LATE for this subgroup is larger. The
increase is consistent across both estimator families and
both covariate specifications.

\paragraph{Kappa estimators: normalised versus scaled variants.}
Among the normalised kappa estimators
($\hat{\tau}^{cb}_{u}$, $\hat{\tau}^{ml}_{u}$,
$\hat{\tau}^{ml}_{a,10}$), estimates cluster between $0.29$
and $0.85$ across the four cells, with the Card specification
yielding the tightest SEs and occasional significance at the
10\% level. These estimates are broadly in line with the
Wald-AIPW results, suggesting the two approaches target a
comparable complier population under the Card specification.

The scaled variants ($\hat{\tau}^{ml}_{a}$,
$\hat{\tau}^{ml}_{t}$, $\hat{\tau}^{ml}_{a,0}$) reveal a
striking pattern under the Kitagawa specification: estimates
inflate substantially---$\hat{\tau}^{ml}_{a,0}$ reaches
$2.712$ for \texttt{educ16}---accompanied by very large
standard errors. This is consistent with near-zero
denominator problems when the Kitagawa propensity score is
sparse, and is further confirmed by the weight diagnostics
(ESS $\approx 1$ for several Kitagawa cells). These
estimates should be interpreted with caution; the weight
diagnostics clearly flag extreme weight concentration that
renders the scaled kappa estimators unreliable in this
specification.



% ============================================================
\section*{Weight Diagnostics}

Table~\ref{tab:wdiag_card_somecol} through
Table~\ref{tab:wdiag_kit_educ16} report four diagnostic
quantities for each estimator and specification cell.
Cells are colour-coded as follows:
\colorbox{diagOK}{\color{diagOKText}green} — Sum\_w $\approx 0$
and ESS $\geq 2$ (no concern);
\colorbox{diagWarn}{\color{diagWarnText}amber} — Sum\_w $\neq 0$
(not translation-invariant);
\colorbox{diagFlag}{\color{diagFlagText}red} — ESS $\leq 1$
(extreme weight concentration).

% helper command: flag a cell red
\newcommand{\cellRed}[1]{%
  \cellcolor{diagFlag}\color{diagFlagText}#1}
\newcommand{\cellAmber}[1]{%
  \cellcolor{diagWarn}\color{diagWarnText}#1}
\newcommand{\cellGreen}[1]{%
  \cellcolor{diagOK}\color{diagOKText}#1}

% ------------------------------------------------------------------
% Table: Card spec, somecol
% ------------------------------------------------------------------
\begin{table}[htbp]
\centering
\begin{threeparttable}
  \caption{Weight diagnostics — Card specification, \texttt{somecol}
           ($D_1 = \mathbf{1}\{\text{educ}>12\}$).}
  \label{tab:wdiag_card_somecol}
  \begin{tabular}{lrrrr}
    \toprule
    Estimator & Sum\_w & ESS & Pct\_neg (\%) & Max\_abs\_w \\
    \midrule
    \multicolumn{5}{l}{\textit{DML estimators}} \\[2pt]
    PLR-IV      & \cellGreen{0.000} & \cellGreen{2} & 33.1 & 0.0327 \\
    Wald-AIPW   & \cellGreen{0.000} & \cellGreen{3} & 31.8 & 0.1456 \\
    \midrule
    \multicolumn{5}{l}{\textit{Kappa estimators}} \\[2pt]
    $\hat\tau^{cb}_u$     & \cellGreen{0.000}  & \cellGreen{3} & 31.8 & 0.0869 \\
    $\hat\tau^{ml}_u$     & \cellGreen{0.000}  & \cellGreen{4} & 31.8 & 0.0519 \\
    $\hat\tau^{ml}_{a,10}$& \cellGreen{0.000}  & \cellGreen{4} & 31.8 & 0.0548 \\
    $\hat\tau^{ml}_a$     & \cellAmber{$-$0.106} & \cellGreen{3} & 31.8 & 0.0545 \\
    $\hat\tau^{ml}_t$     & \cellAmber{$-$0.107} & \cellGreen{3} & 31.8 & 0.0548 \\
    $\hat\tau^{ml}_{a,0}$ & \cellAmber{$-$0.097} & \cellGreen{4} & 31.8 & 0.0495 \\
    \bottomrule
  \end{tabular}
  \begin{tablenotes}[flushleft]\footnotesize
    \item \colorbox{diagOK}{\color{diagOKText}Green}: Sum\_w $=0$, ESS $\geq2$.
          \colorbox{diagWarn}{\color{diagWarnText}Amber}: Sum\_w $\neq 0$ (not translation-invariant).
          \colorbox{diagFlag}{\color{diagFlagText}Red}: ESS $\leq 1$ (extreme weight concentration).
  \end{tablenotes}
\end{threeparttable}
\end{table}

% ------------------------------------------------------------------
% Table: Card spec, educ16
% ------------------------------------------------------------------
\begin{table}[htbp]
\centering
\begin{threeparttable}
  \caption{Weight diagnostics — Card specification, \texttt{educ16}
           ($D_2 = \mathbf{1}\{\text{educ}\geq 16\}$).}
  \label{tab:wdiag_card_educ16}
  \begin{tabular}{lrrrr}
    \toprule
    Estimator & Sum\_w & ESS & Pct\_neg (\%) & Max\_abs\_w \\
    \midrule
    \multicolumn{5}{l}{\textit{DML estimators}} \\[2pt]
    PLR-IV      & \cellGreen{0.000} & \cellRed{1} & 33.1 & 0.0547 \\
    Wald-AIPW   & \cellGreen{0.000} & \cellRed{1} & 31.8 & 0.2765 \\
    \midrule
    \multicolumn{5}{l}{\textit{Kappa estimators}} \\[2pt]
    $\hat\tau^{cb}_u$     & \cellGreen{0.000}  & \cellRed{1} & 31.8 & 0.1971 \\
    $\hat\tau^{ml}_u$     & \cellGreen{0.000}  & \cellRed{1} & 31.8 & 0.0972 \\
    $\hat\tau^{ml}_{a,10}$& \cellGreen{0.000}  & \cellRed{1} & 31.8 & 0.1026 \\
    $\hat\tau^{ml}_a$     & \cellAmber{$-$0.197} & \cellRed{1} & 31.8 & 0.1013 \\
    $\hat\tau^{ml}_t$     & \cellAmber{$-$0.200} & \cellRed{1} & 31.8 & 0.1026 \\
    $\hat\tau^{ml}_{a,0}$ & \cellAmber{$-$0.167} & \cellRed{1} & 31.8 & 0.0855 \\
    \bottomrule
  \end{tabular}
  \begin{tablenotes}[flushleft]\footnotesize
    \item \colorbox{diagOK}{\color{diagOKText}Green}: Sum\_w $=0$, ESS $\geq2$.
          \colorbox{diagWarn}{\color{diagWarnText}Amber}: Sum\_w $\neq 0$ (not translation-invariant).
          \colorbox{diagFlag}{\color{diagFlagText}Red}: ESS $\leq 1$ (extreme weight concentration).
  \end{tablenotes}
\end{threeparttable}
\end{table}

% ------------------------------------------------------------------
% Table: Kitagawa spec, somecol
% ------------------------------------------------------------------
\begin{table}[htbp]
\centering
\begin{threeparttable}
  \caption{Weight diagnostics — Kitagawa specification, \texttt{somecol}
           ($D_1 = \mathbf{1}\{\text{educ}>12\}$).}
  \label{tab:wdiag_kit_somecol}
  \begin{tabular}{lrrrr}
    \toprule
    Estimator & Sum\_w & ESS & Pct\_neg (\%) & Max\_abs\_w \\
    \midrule
    \multicolumn{5}{l}{\textit{DML estimators}} \\[2pt]
    PLR-IV      & \cellGreen{0.000} & \cellGreen{2} & 31.8 & 0.0328 \\
    Wald-AIPW   & \cellGreen{0.000} & \cellGreen{3} & 31.8 & 0.0652 \\
    \midrule
    \multicolumn{5}{l}{\textit{Kappa estimators}} \\[2pt]
    $\hat\tau^{cb}_u$     & \cellGreen{0.000} & \cellGreen{3} & 31.8 & 0.0457 \\
    $\hat\tau^{ml}_u$     & \cellGreen{0.000} & \cellGreen{3} & 31.8 & 0.0444 \\
    $\hat\tau^{ml}_{a,10}$& \cellGreen{0.000} & \cellGreen{3} & 31.8 & 0.0526 \\
    $\hat\tau^{ml}_a$     & \cellAmber{$+$0.305} & \cellGreen{3} & 31.8 & 0.0415 \\
    $\hat\tau^{ml}_t$     & \cellAmber{$+$0.279} & \cellGreen{4} & 31.8 & 0.0379 \\
    $\hat\tau^{ml}_{a,0}$ & \cellAmber{$+$0.387} & \cellGreen{2} & 31.8 & 0.0526 \\
    \bottomrule
  \end{tabular}
  \begin{tablenotes}[flushleft]\footnotesize
    \item \colorbox{diagOK}{\color{diagOKText}Green}: Sum\_w $=0$, ESS $\geq2$.
          \colorbox{diagWarn}{\color{diagWarnText}Amber}: Sum\_w $\neq 0$ (not translation-invariant).
          \colorbox{diagFlag}{\color{diagFlagText}Red}: ESS $\leq 1$ (extreme weight concentration).
  \end{tablenotes}
\end{threeparttable}
\end{table}

% ------------------------------------------------------------------
% Table: Kitagawa spec, educ16
% ------------------------------------------------------------------
\begin{table}[htbp]
\centering
\begin{threeparttable}
  \caption{Weight diagnostics — Kitagawa specification, \texttt{educ16}
           ($D_2 = \mathbf{1}\{\text{educ}\geq 16\}$).}
  \label{tab:wdiag_kit_educ16}
  \begin{tabular}{lrrrr}
    \toprule
    Estimator & Sum\_w & ESS & Pct\_neg (\%) & Max\_abs\_w \\
    \midrule
    \multicolumn{5}{l}{\textit{DML estimators}} \\[2pt]
    PLR-IV      & \cellGreen{0.000} & \cellRed{1} & 31.8 & 0.0570 \\
    Wald-AIPW   & \cellGreen{0.000} & \cellRed{1} & 31.8 & 0.1193 \\
    \midrule
    \multicolumn{5}{l}{\textit{Kappa estimators}} \\[2pt]
    $\hat\tau^{cb}_u$     & \cellGreen{0.000}  & \cellRed{1} & 31.8 & 0.0812 \\
    $\hat\tau^{ml}_u$     & \cellGreen{0.000}  & \cellRed{1} & 31.8 & 0.0783 \\
    $\hat\tau^{ml}_{a,10}$& \cellGreen{0.000}  & \cellRed{0} & 31.8 & 0.1338 \\
    $\hat\tau^{ml}_a$     & \cellAmber{$+$0.586} & \cellRed{1} & 31.8 & 0.0798 \\
    $\hat\tau^{ml}_t$     & \cellAmber{$+$0.496} & \cellRed{1} & 31.8 & 0.0675 \\
    $\hat\tau^{ml}_{a,0}$ & \cellAmber{$+$0.983} & \cellRed{0} & 31.8 & 0.1338 \\
    \bottomrule
  \end{tabular}
  \begin{tablenotes}[flushleft]\footnotesize
    \item \colorbox{diagOK}{\color{diagOKText}Green}: Sum\_w $=0$, ESS $\geq2$.
          \colorbox{diagWarn}{\color{diagWarnText}Amber}: Sum\_w $\neq 0$ (not translation-invariant).
          \colorbox{diagFlag}{\color{diagFlagText}Red}: ESS $\leq 1$ (extreme weight concentration).
  \end{tablenotes}
\end{threeparttable}
\end{table}

% ------------------------------------------------------------------
% Interpretation
% ------------------------------------------------------------------
\subsubsection*{Interpretation of weight diagnostics}

\paragraph{Translation invariance (Sum\_w).}
All DML estimators (PLR-IV and Wald-AIPW) and all three
normalised kappa estimators ($\hat\tau^{cb}_u$,
$\hat\tau^{ml}_u$, $\hat\tau^{ml}_{a,10}$) satisfy
Sum\_w $= 0$ in every cell, confirming translation
invariance: shifting log wages by an additive constant
(e.g.\ switching from dollars to cents) leaves these
estimates unchanged. By contrast, the three scaled
variants ($\hat\tau^{ml}_a$, $\hat\tau^{ml}_t$,
$\hat\tau^{ml}_{a,0}$) have Sum\_w $\neq 0$ in all four
cells without exception. The magnitude grows sharply in
the Kitagawa--\texttt{educ16} cell, where
$\hat\tau^{ml}_{a,0}$ reaches Sum\_w $= 0.983$, nearly
equal to the sum of the positive weights alone. This
explains the extreme sensitivity of $\hat\tau^{ml}_{a,0}$
to outcome rescaling documented in the translation
invariance checks.

\paragraph{Effective sample size (ESS).}
The ESS results draw a sharp boundary along the treatment
dimension. For \texttt{somecol} ($D_1$), all estimators
reach ESS $\geq 2$ in both the Card and Kitagawa
specifications, indicating that the weight mass is
distributed across at least a handful of observations.
For \texttt{educ16} ($D_2$), however, ESS collapses to 1
(or even 0 for $\hat\tau^{ml}_{a,10}$ and
$\hat\tau^{ml}_{a,0}$ under Kitagawa) across every
estimator and both specifications. ESS $= 1$ signals that
a single observation is effectively driving the entire
estimate; ESS $= 0$ arises when the squared weights sum
to more than one, implying the presence of very large
negative weights that inflate the denominator beyond the
sample size.

This pattern is a structural consequence of the rarity
of the \texttt{educ16} treatment: college completion
is much less common in the Card sample than any
post-secondary attendance, yielding a weaker first stage
and more extreme propensity scores near zero or one.
Under these conditions, the kappa weights --- which are
inversely proportional to the instrument propensity score
--- amplify a small number of observations to dominate
the estimate.

\paragraph{Card versus Kitagawa specification.}
Within the \texttt{somecol} subsample the two
specifications produce qualitatively identical diagnostics:
all normalised estimators are clean (green), and the
scaled variants carry the same amber warning in both
cases. The richer Card covariates do not materially
improve the weight distribution here. For
\texttt{educ16}, the diagnostics are uniformly red
under both specifications, confirming that the problem
is driven by the treatment definition rather than by
covariate choice.

\paragraph{Maximum absolute weight.}
For Wald-AIPW under the Card--\texttt{educ16}
specification, Max\_abs\_w reaches $0.277$, by far the
largest value among the DML estimators. Combined with
ESS $= 1$, this signals that a single observation
receives a weight of roughly $28\%$ of the total weight
mass, making the Wald-AIPW estimate in this cell
particularly fragile to individual data points. The
kappa counterpart $\hat\tau^{cb}_u$ has Max\_abs\_w
$= 0.197$ in the same cell, somewhat smaller but still
indicative of the same concentration problem.




% ============================================================
\section*{Learner Comparison: Wald-AIPW with Linear+Logit vs.\ XGBoost}

\subsubsection*{Setup}

We re-estimate the Wald-AIPW estimator using the \texttt{DoubleMLIIVM}
framework of \citet{chernozhukov2018double} with two specific learner choices:
(i) \emph{linear+logit}: OLS regression (\texttt{regr.lm}) for the
outcome nuisance function $g_0$ and logistic regression
(\texttt{classif.log\_reg}) for the instrument and treatment propensity
scores $m_0$, $r_0$; and (ii) \emph{XGBoost}: gradient-boosted trees
(\texttt{regr.xgboost} / \texttt{classif.xgboost}) with
\texttt{alpha=0}, \texttt{subsample=1}, \texttt{max\_delta\_step=0},
\texttt{base\_score=0} and no hyperparameter tuning.
All models use five-fold cross-fitting.

% ------------------------------------------------------------------
% Table: point estimates
% ------------------------------------------------------------------
\begin{table}[htbp]
\centering
\begin{threeparttable}
  \caption{Wald-AIPW point estimates by learner, specification and
           treatment. Outcome: log wages (dollars).}
  \label{tab:learner_estimates}
  \begin{tabular}{llrrrr}
    \toprule
    & & \multicolumn{2}{c}{\texttt{somecol} ($D_1$)}
      & \multicolumn{2}{c}{\texttt{educ16} ($D_2$)} \\
    \cmidrule(lr){3-4}\cmidrule(lr){5-6}
    Learner & & Card & Kitagawa & Card & Kitagawa \\
    \midrule
    \multirow{3}{*}{Linear+logit}
      & Est.    & 0.399 & 0.311 & 0.884 & 0.534 \\
      & (SE)    & (0.248) & (0.235) & (0.608) & (0.409) \\
      & $p$     & 0.108 & 0.186 & 0.146 & 0.192 \\[6pt]
    \multirow{3}{*}{XGBoost}
      & Est.    & \cellRed{$-$13.42} & \cellAmber{$-$0.113}
                & \cellRed{$-$3.198} & \cellRed{$-$3.656} \\
      & (SE)    & \cellRed{(14.41)} & \cellAmber{(0.627)}
                & \cellRed{(4.054)} & \cellRed{(100.9)} \\
      & $p$     & \cellRed{0.352} & \cellAmber{0.857}
                & \cellRed{0.430} & \cellRed{0.971} \\
    \bottomrule
  \end{tabular}
  \begin{tablenotes}[flushleft]\footnotesize
    \item \colorbox{diagFlag}{\color{diagFlagText}Red}:
    economically implausible estimate and/or algebraic identity
    $\boldsymbol{\omega}'\mathbf{Y} = \hat\tau$ fails.
    \colorbox{diagWarn}{\color{diagWarnText}Amber}:
    estimate sign-flipped, algebraic identity fails.
    XGBoost uses \texttt{base\_score=0}, no tuning.
  \end{tablenotes}
\end{threeparttable}
\end{table}

% ------------------------------------------------------------------
% Table: weight diagnostics
% ------------------------------------------------------------------
\begin{table}[htbp]
\centering
\begin{threeparttable}
  \caption{Outcome weight diagnostics for Wald-AIPW by learner.}
  \label{tab:learner_wdiag}
  \setlength{\tabcolsep}{5pt}
  \begin{tabular}{llrrrrc}
    \toprule
    Learner & Spec / Treatment & Sum\_w & ESS & Pct\_neg & Max\_abs\_w
            & Alg.\,check \\
    \midrule
    \multirow{4}{*}{Linear+logit}
      & Card $\cdot$ somecol    & \cellGreen{0.000} & \cellGreen{3}
        & 31.6 & 0.081 & \cellGreen{$\checkmark$} \\
      & Card $\cdot$ educ16     & \cellGreen{0.000} & \cellRed{1}
        & 31.6 & 0.179 & \cellGreen{$\checkmark$} \\
      & Kitagawa $\cdot$ somecol & \cellGreen{0.000} & \cellGreen{3}
        & 31.8 & 0.048 & \cellGreen{$\checkmark$} \\
      & Kitagawa $\cdot$ educ16  & \cellGreen{0.000} & \cellRed{1}
        & 31.8 & 0.082 & \cellGreen{$\checkmark$} \\[4pt]
    \multirow{4}{*}{XGBoost}
      & Card $\cdot$ somecol    & \cellRed{$-$1.281} & \cellRed{0}
        & \cellRed{57.0} & \cellRed{20.03} & \cellRed{$\times$} \\
      & Card $\cdot$ educ16     & \cellRed{$-$0.305} & \cellRed{0}
        & \cellRed{57.0} & \cellRed{4.774} & \cellRed{$\times$} \\
      & Kitagawa $\cdot$ somecol & \cellAmber{$-$0.062} & \cellRed{1}
        & 31.5 & \cellAmber{0.683} & \cellRed{$\times$} \\
      & Kitagawa $\cdot$ educ16  & \cellRed{$-$1.996} & \cellRed{0}
        & 31.5 & \cellRed{22.04} & \cellRed{$\times$} \\
    \bottomrule
  \end{tabular}
  \begin{tablenotes}[flushleft]\footnotesize
    \item Alg.\,check: $\boldsymbol{\omega}'\mathbf{Y} = \hat\tau$
    must hold exactly \citep{knaus2024}.
    ESS $= \bigl(\sum_i\omega_i^2\bigr)^{-1}$.
    ESS\,$= 0$ means $\sum_i\omega_i^2 > 1$, signalling dominant
    negative weights.
    Pct\_neg: percentage of observations with $\omega_i < 0$.
    \colorbox{diagFlag}{\color{diagFlagText}Red}: ESS$\leq 1$
    or $|\text{Sum\_w}|$ large.
    \colorbox{diagWarn}{\color{diagWarnText}Amber}: mild violation.
    \colorbox{diagOK}{\color{diagOKText}Green}: no concern.
  \end{tablenotes}
\end{threeparttable}
\end{table}

% ------------------------------------------------------------------
% Interpretation
% ------------------------------------------------------------------
\subsubsection*{Interpretation}

\paragraph{Linear+logit: consistent and well-behaved.}
The parametric learner delivers estimates that are broadly consistent
with the tuned DML results reported in Table~\ref{tab:dml_estimates}
and with the prior literature on returns to schooling.
Across all four cells the algebraic identity
$\boldsymbol{\omega}'\mathbf{Y} = \hat\tau$ holds exactly, confirming
that the outcome weights are internally coherent and the estimate can
legitimately be expressed as a weighted average of observed log wages.
Sum\_w $= 0$ in all cells, so the estimates are translation-invariant.
The ESS equals 3 for the \texttt{somecol} cells and collapses to 1 for
\texttt{educ16}, consistent with the weight concentration documented in
Table~\ref{tab:wdiag_card_educ16} and \ref{tab:wdiag_kit_educ16} --- a
feature of the treatment rarity rather than the learner choice.

\paragraph{XGBoost: catastrophic failure of the algebraic identity.}
The XGBoost specification produces economically nonsensical estimates in
all four cells: $-13.42$ (Card/\texttt{somecol}), $-3.20$
(Card/\texttt{educ16}), $-0.11$ (Kitagawa/\texttt{somecol}), and
$-3.66$ with a standard error of $100.9$
(Kitagawa/\texttt{educ16}). All four fail the algebraic identity
$\boldsymbol{\omega}'\mathbf{Y} = \hat\tau$, meaning the outcome weights
do not correspond to the actual estimate --- a fundamental structural
breakdown, not statistical noise.

\paragraph{Why XGBoost fails: the \texttt{base\_score=0} misconfiguration.}
\citet{knaus2024treatment} establishes that availability and properties of outcome
weights depend critically on implementation choices (his Table~5 and~6).
Specifically, he shows that \emph{Wald-AIPW with XGBoost fails
Condition~C3} --- the smoothness condition required for the outcome
weights to be normalised --- and documents in his empirical Monte Carlo
simulation (EMCS) that outliers ranging from $-16$ to $55$ arise in
exactly this learner--estimator combination
(\citealt{knaus2024treatment}, Figure~1 and Table~6).
The root cause in our application is the XGBoost parameter
\texttt{base\_score=0}. In XGBoost $\geq 1.6$ (the version available on
CRAN at the time of writing) setting \texttt{base\_score=0} no longer
produces a zero-mean base learner; instead it biases all initial leaf
predictions severely towards zero, distorting the residuals that drive
subsequent tree splits \citep{knaus2024treatment}.
In earlier versions, \texttt{base\_score=0} initialised every
observation to zero and made the smoother matrix \emph{affine} --- the
precise requirement for $\boldsymbol{\omega}'\mathbf{Y} = \hat\tau$ to
hold. Condition~C3 fails under the current implementation because the
cross-fitting folds produce propensity score predictions saturated near
zero or one, generating enormous AIPW correction terms when inverted.
This manifests in ESS $= 0$ (i.e.\ $\sum_i\omega_i^2 > 1$),
Max$|\omega_i| = 20$ for Card/\texttt{somecol} and $22$ for
Kitagawa/\texttt{educ16}, and Pct\_neg rising to 57\% for the Card
specification compared to $\approx 32\%$ everywhere else.

Without hyperparameter tuning the gradient-boosted trees overfit the
nuisance functions in-sample; cross-fitting partially mitigates this but
cannot recover the affine smoother property that the algebraic identity
requires. 





% ============================================================
% ============================================================
\section*{Diagnostic Evidence: Why the Results Look the Way They Do}

The following tables collect the key structural diagnostics from the
Card (1995) sample that explain the patterns documented in all preceding
sections.  The diagnostics fall into four groups: first-stage strength,
compliance heterogeneity, propensity score distributions, and the
profile of extreme-weight observations.

% ------------------------------------------------------------------
% Table: first stage + compliance
% ------------------------------------------------------------------
\begin{table}[htbp]
\centering
\begin{threeparttable}
  \caption{First-stage strength and compliance type decomposition.
           Instrument: \texttt{nearc4}. Sample: $N = 3{,}010$.}
  \label{tab:firststage}
  \begin{tabular}{lrrrrrr}
    \toprule
    & \multicolumn{3}{c}{First stage}
    & \multicolumn{3}{c}{Compliance shares} \\
    \cmidrule(lr){2-4}\cmidrule(lr){5-7}
    Treatment & $\hat\pi_\text{fs}$ & $F$ (no $X$) & $F$ (Card $X$)
              & $\hat\pi_c$ & $\hat\pi_\text{at}$ & $\hat\pi_\text{nt}$ \\
    \midrule
    \texttt{somecol} ($D_1$) & \cellGreen{0.122} & \cellGreen{39.3}
      & \cellAmber{$\approx 11.8$}
      & \cellAmber{12.2\%} & 42.2\% & 45.6\% \\
    \texttt{educ16} ($D_2$)  & \cellRed{0.069}   & \cellAmber{15.6}
      & \cellRed{$\approx 3.2$}
      & \cellRed{6.9\%}   & 22.5\% & 70.7\% \\
    \bottomrule
  \end{tabular}
  \begin{tablenotes}[flushleft]\footnotesize
    \item $\hat\pi_\text{fs} = \Pr(D=1\mid Z=1) - \Pr(D=1\mid Z=0)$.
    $F$ (Card $X$) is the partial $F$-statistic on $Z$ in a first-stage
    OLS with the full Card covariate set.
    $\hat\pi_c$, $\hat\pi_\text{at}$, $\hat\pi_\text{nt}$: estimated
    complier, always-taker, and never-taker shares.
  \end{tablenotes}
\end{threeparttable}
\end{table}

% ------------------------------------------------------------------
% Table: propensity score distributions
% ------------------------------------------------------------------
\begin{table}[htbp]
\centering
\begin{threeparttable}
  \caption{Instrument propensity score $\hat{p}(Z{=}1\mid X)$ and
           treatment propensity $\hat{P}(D{=}1\mid X)$ distributions
           (Card specification).}
  \label{tab:propscores}
  \setlength{\tabcolsep}{5pt}
  \begin{tabular}{lrrrrrrrr}
    \toprule
    & Min & Q1 & Median & Mean & Q3 & Max
    & $\%< 0.3$ & $\%> 0.7$ \\
    \midrule
    \multicolumn{9}{l}{\textit{Instrument propensity $\hat{p}(Z{=}1\mid X)$}} \\[2pt]
    Card spec      & 0.190 & 0.430 & \cellAmber{0.791} & 0.682 & 0.874
                   & 0.950 & \cellAmber{7.9\%} & \cellRed{63.9\%} \\
    Kitagawa spec  & 0.253 & 0.455 & \cellAmber{0.798} & 0.682 & 0.878
                   & 0.907 & 7.3\% & \cellRed{63.8\%} \\[4pt]
    \multicolumn{9}{l}{\textit{Treatment propensity $\hat{P}(D{=}1\mid X)$,
                         Card spec}} \\[2pt]
    \texttt{somecol} ($D_1$) & 0.180 & 0.393 & 0.577 & 0.505 & 0.596
                              & 0.812 & \cellGreen{0\%} & \cellGreen{0\%} \\
    \texttt{educ16} ($D_2$)  & \cellRed{0.067} & \cellFlag{0.214}
                              & \cellFlag{0.302} & \cellAmber{0.271}
                              & 0.347 & \cellRed{0.440}
                              & \cellGreen{0\%} & \cellGreen{0\%} \\
    \bottomrule
  \end{tabular}
  \begin{tablenotes}[flushleft]\footnotesize
    \item Last two columns give the share of observations below 0.3 and
    above 0.7 respectively.
    Both instrument propensity distributions are highly right-skewed
    (median $\approx 0.79$); the Card regional dummies push some
    $\hat{p}_i$ toward 0.19 (rural South) or 0.95 (North-East cities).
    The \texttt{educ16} treatment propensity tops out at 0.44 --- no
    observation is near-certain to complete college given observables.
  \end{tablenotes}
\end{threeparttable}
\end{table}

% ------------------------------------------------------------------
% Table: kappa weight profile
% ------------------------------------------------------------------
\begin{table}[htbp]
\centering
\begin{threeparttable}
  \caption{Kappa weight distribution and extreme-observation
           concentration (Card specification).}
  \label{tab:kappa_profile}
  \begin{tabular}{lrrrr}
    \toprule
    & \multicolumn{2}{c}{\texttt{somecol} ($D_1$)}
    & \multicolumn{2}{c}{\texttt{educ16} ($D_2$)} \\
    \cmidrule(lr){2-3}\cmidrule(lr){4-5}
    Metric & Value & Note & Value & Note \\
    \midrule
    Min $\kappa_i$         & \cellRed{$-$13.87} & &
                             \cellRed{$-$13.87} & \\
    Median $\kappa_i$      & $+$1.000 & complier & \cellRed{$-$0.072}
                           & \textit{negative!} \\
    Share $\kappa_i < 0$   & \cellAmber{44.5\%} & &
                             \cellRed{55.4\%} & \\
    Top-5 share of $|\kappa|$  & \cellRed{2.3\%} & &
                                  \cellRed{2.7\%} & \\
    Top-20 share of $|\kappa|$ & \cellRed{7.1\%} & &
                                   \cellRed{6.3\%} & \\
    \midrule
    \multicolumn{5}{l}{\textit{Profile of top extreme observations}} \\[2pt]
    $\hat{p}(Z{=}1\mid X)$ & \multicolumn{2}{c}{0.933} &
                              \multicolumn{2}{c}{0.933} \\
    Observed $Z$           & \multicolumn{2}{c}{0} &
                              \multicolumn{2}{c}{0} \\
    Observed $D$           & \multicolumn{2}{c}{1} &
                              \multicolumn{2}{c}{1} \\
    Implied $\kappa$       & \multicolumn{4}{c}
      {$\frac{0}{0.933} - \frac{1}{1-0.933} \approx -13.9$} \\
    \bottomrule
  \end{tabular}
  \begin{tablenotes}[flushleft]\footnotesize
    \item The extreme observations are \emph{``surprising
    non-recipients''}: observations where all covariates predict a very
    high probability of receiving the instrument ($\hat{p} \approx
    0.93$) but the instrument was in fact not received ($Z = 0$), and
    who are nonetheless treated ($D = 1$). Their kappa weight formula
    $\kappa_i = Z_i/\hat{p}_i - (1-Z_i)/(1-\hat{p}_i)$ evaluates to
    $0/0.933 - 1/0.067 \approx -13.9$.
  \end{tablenotes}
\end{threeparttable}
\end{table}

% ------------------------------------------------------------------
% Table: risk-group audit
% ------------------------------------------------------------------
\begin{table}[htbp]
\centering
\begin{threeparttable}
  \caption{Risk-group audit: high-$\hat{p}(Z{=}1\mid X)$ observations
           with $Z = 0$, $D = 1$ (Card specification).}
  \label{tab:riskgroup}
  \begin{tabular}{lrrrr}
    \toprule
    $\hat{p}$ cutoff & $N$ at risk & Share of sample
                     & Mean $\kappa$ & Max$\,|\kappa|$ \\
    \midrule
    \multicolumn{5}{l}{\textit{Treatment: \texttt{somecol} ($D_1$)}} \\[2pt]
    $\hat{p} > 0.70$ & 139 & \cellAmber{4.6\%} &
                       \cellRed{$-$4.85} & \cellRed{13.87} \\
    $\hat{p} > 0.80$ &  59 & 2.0\% &
                       \cellRed{$-$7.11} & \cellRed{13.87} \\
    $\hat{p} > 0.90$ &  11 & 0.4\% &
                       \cellRed{$-$13.58} & \cellRed{13.87} \\[4pt]
    \multicolumn{5}{l}{\textit{Treatment: \texttt{educ16} ($D_2$)}} \\[2pt]
    $\hat{p} > 0.70$ &  78 & \cellAmber{2.6\%} &
                       \cellRed{$-$4.54} & \cellRed{13.87} \\
    $\hat{p} > 0.80$ &  29 & 1.0\% &
                       \cellRed{$-$6.95} & \cellRed{13.87} \\
    $\hat{p} > 0.90$ &   5 & 0.2\% &
                       \cellRed{$-$13.55} & \cellRed{13.87} \\
    \bottomrule
  \end{tabular}
  \begin{tablenotes}[flushleft]\footnotesize
    \item ``At risk'' = $\hat{p}(Z{=}1\mid X)$ above threshold, $Z=0$,
    $D=1$.  The 11 observations (5 for \texttt{educ16}) with
    $\hat{p}>0.90$ carry mean kappa of $-13.6$ and are responsible for
    a disproportionate share of total weight mass, driving the ESS to 1.
  \end{tablenotes}
\end{threeparttable}
\end{table}

% ------------------------------------------------------------------
% Interpretation
% ------------------------------------------------------------------
\subsubsection*{How the diagnostics explain the estimation results}

\paragraph{Weak first stage for \texttt{educ16}.}
Once the full Card covariate set is included, the partial $F$-statistic
on \texttt{nearc4} in the \texttt{educ16} first stage falls to
approximately 3.2 (Table~\ref{tab:firststage}), well below the
conventional weak-instrument threshold of 10.  The first-stage
difference $\hat\pi_\text{fs} = 0.069$ means that the Wald numerator is
scaled by a very small quantity, inflating the LATE estimate and its
variance.  This propagates directly into the ESS collapse observed in
Tables~\ref{tab:wdiag_card_educ16} and~\ref{tab:wdiag_kit_educ16}:
the outcome weights must compensate for the weak signal by
concentrating on the small number of observations where the instrument
variation is informative.

\paragraph{Extreme non-complier majority.}
Only 6.9\% of the sample are compliers for \texttt{educ16}
(Table~\ref{tab:firststage}).  By construction, only compliers can
receive positive kappa weights; all always-takers and never-takers
receive negative weights.  The resulting 55\% negative-kappa share for
\texttt{educ16} (Table~\ref{tab:kappa_profile}) is therefore not a
numerical artefact but a direct consequence of the compliance structure.
This also explains why the median kappa for \texttt{educ16} is negative
($-0.072$): more than half the sample sits on the wrong side of the
weighting scheme.

\paragraph{Right-skewed instrument propensity and the ``surprising
non-recipient'' mechanism.}
The instrument propensity $\hat{p}(Z{=}1\mid X)$ is heavily
right-skewed with median $\approx 0.79$ in both specifications
(Table~\ref{tab:propscores}).  Observations with $\hat{p}_i \approx
0.93$ but $Z_i = 0$ and $D_i = 1$ receive kappa weights of
$\approx -13.9$ (Table~\ref{tab:kappa_profile}): they are ``surprising
non-recipients'' who, by all observable characteristics, should have
had a college nearby but did not, yet attended college regardless.  Only
11 such observations (5 for \texttt{educ16}) lie above the
$\hat{p}>0.90$ threshold, but they carry mean kappa of $-13.6$
(Table~\ref{tab:riskgroup}).  When any machine learning nuisance model
--- in particular the untuned XGBoost --- assigns these observations
extreme propensity scores, the AIPW correction terms diverge, the
algebraic identity $\boldsymbol{\omega}'\mathbf{Y}=\hat\tau$ fails, and
the estimates collapse into the $(-13, -3)$ range reported in
Table~\ref{tab:learner_estimates}.

\paragraph{Why the Card spec does better than Kitagawa.}
The Card covariate set includes regional dummies that successfully
separate the few observations with genuinely low instrument propensity
($\hat{p} < 0.19$, rural South) from the majority.  The Kitagawa spec
lacks these dummies, so its propensity model mixes observations from
different regions into a single intermediate propensity band, producing
a lower implied ESS (2.9 vs.\ 4.6, Table~\ref{tab:propscores}) even
though the raw propensity distributions look similar.  More covariates
allow the model to locate the complier subpopulation more precisely,
which is why covariate richness matters for this observational IV
design, in contrast to near-randomised settings where specification
choice has little effect on the weight diagnostics.





% Here is where the citations start 
\bibliographystyle{apalike}
\bibliography{references}
\end{document}